
\begin{UPSTIexercice}{Créer une structure cohérente}
    On regroupe les données dans une seule entité : un \textbf{type structuré}.
    
    \UPSTIquestion{Compléter la définition du type \texttt{T\_Etudiant} (prénom, âge, moyenne)}
    \begin{lstlisting}[language=c,escapeinside={(*}{*)},caption={T\_Etudiant, à compléter}]
typedef struct {
    (*\trou[string prenom;]{7.5cm}*)
    (*\trou[int age;]{7.5cm}*)
    (*\trou[float moyenne;]{7.5cm}*)
} T_Etudiant;
\end{lstlisting}
    \UPSTIquestion{Ecrire les les déclarations des variables de type \texttt{T\_Etudiant} suivantes :}
    \begin{itemize}
        \item une variable \texttt{etudiant\_1}, puis sur une seconde ligne, nommer cet étudiant "Bob", lui donner 22 ans et une moyenne de 15.5,
        \item Une variable \texttt{etudiant\_2} qui sera initialisée (sur la ligne de déclaration) directement avec les valeurs "Alice", 20 et 17.0,
        \item un tableau \texttt{classe} de 30 étudiants.
    \end{itemize}



\begin{lstlisting}[language=c,escapeinside={(*}{*)},caption={Saisie et recherche du meilleur étudiant, à compléter}]
    // Etudiant 1 : 
    T_Etudiant (*\trou[etudiant\_1;]{7.5cm}*); 
    (*\trou[etudiant\_1.prenom = "Bob";]{7.5cm}*)
    (*\trou[etudiant\_1.age = 22;]{7.5cm}*)
    (*\trou[etudiant\_1.moyenne = 15.5f;]{7.5cm}*)
    // Etudiant 2 : 
    (*\trou[etudiant\_2 = \{"Alice", 20, 17.0f\};]{13.5cm}*); 
\end{lstlisting}
\UPSTIquestion{Ecrire la déclaration et la saisie d'un tableau de 30 étudiants, puis la recherche du meilleur étudiant (celui ayant la meilleure moyenne).}
\begin{lstlisting}[language=c,escapeinside={(*}{*)},caption={Saisie et recherche du meilleur étudiant, à compléter}]
    // Classe : 
    (*\trou[{T\_Etudiant classe[30];}]{9.5cm}*) 

for (inti = 0; i < 30; i++) {
    printf("Etudiant %d\n", i + 1);
    (*\trou[{classe[i].prenom}]{7.5cm}*) = get_string("Prenom : ");
    (*\trou[{classe[i].age}]{7.5cm}*) = get_int("Age : ");
    (*\trou[{classe[i].moyenne}]{7.5cm}*) = get_float("Moyenne : ");
}

int indice_meilleur = 0;
for (i = 1; i < 30; i++) {
    if ((*\trou[{classe[i].moyenne > classe[indice\_meilleur].moyenne}]{11.5cm}*)) {
        indice_meilleur = i;
    }
}

printf("Meilleur etudiant : %s, moyenne = %.2f\n",
       (*\trou[{classe[indice\_meilleur].prenom}]{7.5cm}*),
       (*\trou[{classe[indice\_meilleur].moyenne}]{7.5cm}*));
\end{lstlisting}
\end{UPSTIexercice}

% ==========================================================================

\begin{UPSTIexercice}{Représenter une heure}
    On veut manipuler une heure sous la forme \texttt{HH:MM:SS}.

    \UPSTIquestion{Compléter le type \texttt{T\_Heure}.}

    \begin{lstlisting}[language=c,escapeinside={(*}{*)},caption={T\_Heure, à compléter}]
typedef struct {
    (*\trou[int h;]{6.5cm}*)
    (*\trou[int m;]{6.5cm}*)
    (*\trou[int s;]{6.5cm}*)
} T_Heure;
\end{lstlisting}

    \UPSTIquestion{Écrire le corps des fonctions permettant d'afficher une heure
        (\texttt{afficher\_heure}), de la convertir en secondes (\texttt{en\_secondes}), et
        inversement de reconstruire une heure à partir d'un nombre de secondes (\texttt{inverse}).}

    \begin{lstlisting}[language=c,caption={Prototypes}]
void afficher_heure(T_Heure t);
int en_secondes(T_Heure t);
T_Heure inverse(int total_secondes);
\end{lstlisting}

\begin{UPSTIeleveOnlyEnv}
    \UPSTIquadrillage{7}
\end{UPSTIeleveOnlyEnv}

\begin{UPSTIcorrectionP}{Fonctions sur T\_Heure}
\begin{lstlisting}[language=c,caption={Correction}]
void afficher_heure(T_Heure t) {
    printf("%02d:%02d:%02d\n", t.h, t.m, t.s);
}

int en_secondes(T_Heure t) {
    return t.h * 3600 + t.m * 60 + t.s;
}

T_Heure inverse(int total_secondes) {
    T_Heure t;
    t.h = total_secondes / 3600;
    total_secondes %= 3600;
    t.m = total_secondes / 60;
    t.s = total_secondes % 60;
    return t;
}
\end{lstlisting}
\end{UPSTIcorrectionP}

    \begin{lstlisting}[language=c,escapeinside={(*}{*)},caption={Utilisation dans le main, à compléter}]
int main(void) {
    T_Heure h1 = {8, 30, 0};
    T_Heure h2 = {10, 15, 30};

    int s1 = (*\trou[en\_secondes(h1)]{7.5cm}*);
    int s2 = (*\trou[en\_secondes(h2)]{7.5cm}*);

    printf("Difference = %d secondes\n", s2 - s1);
    return 0;
}
\end{lstlisting}
\end{UPSTIexercice}

% --------------------------------------------------------------------------

\begin{UPSTIexercice}{Gérer un produit en stock}
    On veut garder les informations suivantes pour chaque produit :
    \begin{itemize}
        \item son nom (chaîne de caractères),
        \item son prix HT (nombre à virgule),
        \item son taux de TVA (nombre à virgule).
    \end{itemize}

    \UPSTIquestion{Compléter le type \texttt{T\_Produit}.}

    \begin{lstlisting}[language=c,escapeinside={(*}{*)},caption={T\_Produit, à compléter}]
typedef struct {
    (*\trou[string nom;]{7.5cm}*)
    (*\trou[float prixHT;]{7.5cm}*)
    (*\trou[float tva;]{7.5cm}*)  // en pourcentage, par ex. 20.0
} T_Produit;
\end{lstlisting}

    \UPSTIquestion{Écrire le corps de la fonction renvoyant le prix TTC d'un produit (la fonction
        ne prendra qu'un seul argument).}

    \begin{lstlisting}[language=c,caption={Prototype}]
float prixTTC(T_Produit p);
\end{lstlisting}

\begin{UPSTIeleveOnlyEnv}
    \UPSTIquadrillage{3}
\end{UPSTIeleveOnlyEnv}

\begin{UPSTIcorrectionP}{prixTTC}
\begin{lstlisting}[language=c,caption={Correction}]
float prixTTC(T_Produit p) {
    return p.prixHT * (1.0f + p.tva / 100.0f);
}
\end{lstlisting}
\end{UPSTIcorrectionP}
\end{UPSTIexercice}

% --------------------------------------------------------------------------

\begin{UPSTIinfor}{La fonction racine carrée}
    Le calcul d'une norme nécessite une racine carrée. Le langage C fournit cette fonction dans
    la bibliothèque \texttt{math.h}, qu'il faut inclure avec \texttt{\#include <math.h>} :
    \begin{itemize}
        \item \texttt{double sqrt(double x)} renvoie $\sqrt{x}$ ;
        \item \texttt{float sqrtf(float x)} fait de même pour un argument et un résultat en \texttt{float}.
    \end{itemize}
    On utilisera \texttt{sqrtf} lorsque les champs de la structure sont en \texttt{float}, afin d'éviter
    les conversions implicites \texttt{float}/\texttt{double}.
\end{UPSTIinfor}

% --------------------------------------------------------------------------

\begin{UPSTIexercice}{Gestion d’un vecteur 2D}
    On manipule des vecteurs dans un plan :
    \begin{lstlisting}[language=C]
typedef struct{
    float x;
    float y;
} T_Vecteur;
\end{lstlisting}

    \UPSTIquestion{Écrire le corps des fonctions \texttt{norme}, \texttt{addition} et \texttt{oppose}.}

    \begin{lstlisting}[language=c,caption={Prototypes}]
#include <math.h>

float norme(T_Vecteur v);
T_Vecteur addition(T_Vecteur a, T_Vecteur b);
T_Vecteur oppose(T_Vecteur v);
\end{lstlisting}

\begin{UPSTIeleveOnlyEnv}
    \UPSTIquadrillage{6}
\end{UPSTIeleveOnlyEnv}

\begin{UPSTIcorrectionP}{Fonctions sur T\_Vecteur}
\begin{lstlisting}[language=c,caption={Correction}]
float norme(T_Vecteur v) {
    return sqrtf(v.x * v.x + v.y * v.y);
}

T_Vecteur addition(T_Vecteur a, T_Vecteur b) {
    T_Vecteur res;
    res.x = a.x + b.x;
    res.y = a.y + b.y;
    return res;
}

T_Vecteur oppose(T_Vecteur v) {
    T_Vecteur res;
    res.x = -v.x;
    res.y = -v.y;
    return res;
}
\end{lstlisting}
\end{UPSTIcorrectionP}

    \begin{lstlisting}[language=c,escapeinside={(*}{*)},caption={Utilisation dans le main, à compléter}]
int main(void) {
    T_Vecteur a = {3.0f, 4.0f};
    T_Vecteur b = {1.0f, 2.0f};

    printf("||a|| = %f\n", (*\trou[norme(a)]{7.5cm}*));
    printf("||b|| = %f\n", (*\trou[norme(b)]{7.5cm}*));

    T_Vecteur diff = (*\trou[{addition(a, oppose(b))}]{11.5cm}*);
    printf("||a-b|| = %f\n", norme(diff));
    return 0;
}
\end{lstlisting}
\end{UPSTIexercice}

% ==========================================================================

\begin{UPSTIexercice}{Extension : struct imbriquées}
    Il est également possible de créer des structures imbriquées, c'est-à-dire d'utiliser une structure comme champ d'une autre structure. Par exemple, on peut définir une structure \texttt{Adresse} et l'utiliser dans une structure \texttt{Etudiant} :
    \begin{lstlisting}[language=C]
typedef struct {
    string rue;
    int codePostal;
    string ville;
} T_Adresse;

typedef struct {
    string prenom;
    int age;
    float moyenne;
    T_Adresse adresse;
} T_EtudiantAdresse;
\end{lstlisting}

    \UPSTIquestion{Quel est l'avantage de créer une structure imbriquée plutôt qu'ajouter 3 champs à \texttt{Etudiant} ?}
    \begin{UPSTIeleveOnlyEnv}
        \UPSTIquadrillage{3}
    \end{UPSTIeleveOnlyEnv}
    \begin{UPSTIcorrectionP}{Avantage des structures imbriquées}
        L'avantage de créer une structure imbriquée est de pouvoir réutiliser T\_Adresse pour d'autres structures (Professeur, Entreprise, etc.) et de faciliter les modifications ultérieures.
    \end{UPSTIcorrectionP}

    \UPSTIquestion{Compléter l'accès à la ville d'un étudiant.}

    \begin{lstlisting}[language=c,escapeinside={(*}{*)},caption={Accès, à compléter}]
// Si e est une variable de type T_EtudiantAdresse :
string ville_e = (*\trou[{e.adresse.ville}]{7.5cm}*);
\end{lstlisting}

    \UPSTIquestion{Écrire le corps de la fonction \texttt{habite\_ville}.}

    \begin{lstlisting}[language=c,caption={Prototype}]
int habite_ville(T_EtudiantAdresse e, const char* ville);
\end{lstlisting}

\begin{UPSTIeleveOnlyEnv}
    \UPSTIquadrillage{2.5}
\end{UPSTIeleveOnlyEnv}

\begin{UPSTIcorrectionP}{habite\_ville}
\begin{lstlisting}[language=c,caption={Correction}]
int habite_ville(T_EtudiantAdresse e, const char* ville) {
    return strcmp(e.adresse.ville, ville) == 0;
}
\end{lstlisting}
\end{UPSTIcorrectionP}
\end{UPSTIexercice}

% ==========================================================================

\begin{UPSTIinfor}{Les types énumérés}
    Un type \textbf{énuméré} permet de nommer un ensemble fini de valeurs possibles, plutôt que
    d'utiliser des entiers « magiques » dont la signification n'est pas claire. Il se définit avec
    le mot-clé \texttt{enum} :
    \begin{lstlisting}[language=C]
typedef enum {
    VALEUR1,
    VALEUR2,
    VALEUR3
} T_NomDuType;
\end{lstlisting}
    Par défaut, \texttt{VALEUR1} vaut \texttt{0}, \texttt{VALEUR2} vaut \texttt{1}, etc. : un type énuméré
    reste manipulé comme un entier, mais le nom choisi rend le code bien plus lisible qu'un simple
    \texttt{0}, \texttt{1} ou \texttt{2}.
\end{UPSTIinfor}

\begin{UPSTIexercice}{Niveau d'un étudiant}
    On souhaite associer à chaque étudiant un \textbf{niveau} parmi \texttt{L1}, \texttt{L2} et \texttt{L3},
    et stocker ce niveau dans la structure \texttt{T\_EtudiantAdresse} définie précédemment.

    \UPSTIquestion{Compléter la définition du type énuméré \texttt{T\_Niveau}, puis l'ajout du champ
        \texttt{niveau} dans \texttt{T\_EtudiantAdresse}.}

    \begin{lstlisting}[language=c,escapeinside={(*}{*)},caption={T\_Niveau, à compléter}]
typedef enum {
    (*\trou[L1,]{6.5cm}*)
    (*\trou[L2,]{6.5cm}*)
    (*\trou[L3]{6.5cm}*)
} T_Niveau;

typedef struct {
    string prenom;
    int age;
    float moyenne;
    T_Adresse adresse;
    (*\trou[{T\_Niveau niveau;}]{7.5cm}*)
} T_EtudiantAdresse;
\end{lstlisting}

    \UPSTIquestion{Compléter la fonction \texttt{afficher\_niveau} qui affiche en toutes lettres
        (« L1 », « L2 » ou « L3 ») le niveau d'un étudiant, à l'aide d'un \texttt{switch}.}

    \begin{lstlisting}[language=c,escapeinside={(*}{*)},caption={afficher\_niveau, à compléter}]
void afficher_niveau(T_EtudiantAdresse e) {
    switch (e.niveau) {
        case L1:
            (*\trou[printf("L1\textbackslash n");]{7.5cm}*)
            break;
        case L2:
            (*\trou[printf("L2\textbackslash n");]{7.5cm}*)
            break;
        case L3:
            (*\trou[printf("L3\textbackslash n");]{7.5cm}*)
            break;
    }
}
\end{lstlisting}

    \UPSTIquestion{Écrire le corps de la fonction \texttt{compter\_niveau} qui compte dans un
        tableau \texttt{promo} le nombre d'étudiants d'un niveau donné.}

    \begin{lstlisting}[language=c,caption={Prototype}]
int compter_niveau(T_EtudiantAdresse promo[], int nbEtudiants, T_Niveau niveau);
\end{lstlisting}

\begin{UPSTIeleveOnlyEnv}
    \UPSTIquadrillage{4}
\end{UPSTIeleveOnlyEnv}

\begin{UPSTIcorrectionP}{compter\_niveau}
\begin{lstlisting}[language=c,caption={Correction}]
int compter_niveau(T_EtudiantAdresse promo[], int nbEtudiants, T_Niveau niveau) {
    int i, compteur = 0;

    for (i = 0; i < nbEtudiants; i++) {
        if (promo[i].niveau == niveau) {
            compteur++;
        }
    }
    return compteur;
}
\end{lstlisting}
\end{UPSTIcorrectionP}
\end{UPSTIexercice}

% ==========================================================================

\begin{UPSTIexercice}{Pour aller plus loin : tableau d'étudiants avec adresse}
    On dispose maintenant d'une \textbf{promotion entière} d'étudiants, chacun avec son adresse
    (type \texttt{T\_EtudiantAdresse} défini précédemment). On fixe une taille maximale de promotion :
    \begin{lstlisting}[language=C]
#define TAILLE_MAX_PROMO 40
\end{lstlisting}

    \UPSTIquestion{Compléter la déclaration du tableau de la promotion, et la saisie de ses
        \texttt{nbEtudiants} étudiants (prénom, âge, moyenne, puis les 3 champs de l'adresse).}

    \begin{lstlisting}[language=c,escapeinside={(*}{*)},caption={Déclaration et saisie, à compléter}]
(*\trou[{T\_EtudiantAdresse promo[TAILLE\_MAX\_PROMO];}]{13.5cm}*)
int nbEtudiants;
int i;

nbEtudiants = get_int("Nombre d'etudiants : ");

for (i = 0; i < nbEtudiants; i++) {
    promo[i].prenom = get_string("Prenom : ");
    promo[i].age    = get_int("Age : ");
    promo[i].moyenne = get_float("Moyenne : ");
    (*\trou[{promo[i].adresse.rue}]{7.5cm}*)        = get_string("Rue : ");
    (*\trou[{promo[i].adresse.codePostal}]{7.5cm}*) = get_int("Code postal : ");
    (*\trou[{promo[i].adresse.ville}]{7.5cm}*)      = get_string("Ville : ");
}
\end{lstlisting}

    \UPSTIquestion{Sur le même modèle que \texttt{compter\_niveau}, et en réutilisant
        \texttt{habite\_ville}, compléter le corps de la fonction qui compte combien d'étudiants
        habitent une ville donnée.}

    \begin{lstlisting}[language=c,escapeinside={(*}{*)},caption={compter\_ville, à compléter}]
int compter_ville(T_EtudiantAdresse promo[], int nbEtudiants, const char *ville) {
    int i, compteur = 0;

    for (i = 0; i < nbEtudiants; i++) {
        if ((*\trou[{habite\_ville(promo[i], ville)}]{9.5cm}*)) {
            compteur++;
        }
    }
    return compteur;
}
\end{lstlisting}

    \UPSTIquestion{Pourquoi la fonction \texttt{compter\_ville} doit-elle recevoir \texttt{nbEtudiants} en paramètre, en plus du tableau \texttt{promo} ?}
    \begin{UPSTIeleveOnlyEnv}
        \UPSTIquadrillage{3}
    \end{UPSTIeleveOnlyEnv}
    \begin{UPSTIcorrectionP}{Pourquoi nbEtudiants est nécessaire}
        La fonction \texttt{compter\_ville} a besoin de connaître le nombre d'étudiants pour pouvoir itérer sur le tableau \texttt{promo} et compter les étudiants qui habitent la ville donnée.
    \end{UPSTIcorrectionP}
\end{UPSTIexercice}

% ==========================================================================

\begin{UPSTIexercice}{Pour aller plus loin : filtrer et afficher par ville}
    On souhaite maintenant afficher la \textbf{liste complète} des étudiants d'une ville donnée
    (prénom et moyenne), puis calculer la \textbf{moyenne de la classe restreinte aux étudiants
    d'une ville donnée}. On réutilise le tableau \texttt{promo} et la fonction \texttt{compter\_ville}
    de l'exercice précédent.

    \UPSTIquestion{En réutilisant \texttt{habite\_ville} et \texttt{compter\_ville}, compléter le
        corps de la fonction \texttt{afficher\_ville} qui affiche le prénom et la moyenne de
        chaque étudiant d'une ville donnée, et de la fonction \texttt{moyenne\_ville} qui renvoie
        la moyenne de ces étudiants.}

    \begin{lstlisting}[language=c,escapeinside={(*}{*)},caption={afficher\_ville et moyenne\_ville, à compléter}]
void afficher_ville(T_EtudiantAdresse promo[], int nbEtudiants, const char *ville) {
    int i;

    for (i = 0; i < nbEtudiants; i++) {
        if ((*\trou[{habite\_ville(promo[i], ville)}]{9.5cm}*)) {
            printf("%s : %.2f\n", promo[i].prenom, promo[i].moyenne);
        }
    }
}

float moyenne_ville(T_EtudiantAdresse promo[], int nbEtudiants, const char *ville) {
    int i;
    float somme = 0.0f;
    int nb = (*\trou[{compter\_ville(promo, nbEtudiants, ville)}]{11.5cm}*);

    for (i = 0; i < nbEtudiants; i++) {
        if (habite_ville(promo[i], ville)) {
            somme += promo[i].moyenne;
        }
    }

    if (nb == 0) {
        return 0.0f;
    }
    return somme / nb;
}
\end{lstlisting}

    \UPSTIquestion{Dans le \texttt{main}, écrire l'appel permettant d'afficher les étudiants de la ville saisie par l'utilisateur, puis leur moyenne.}

    \begin{lstlisting}[language=c,escapeinside={(*}{*)},caption={Utilisation dans le main, à compléter}]
int main(void) {
    T_EtudiantAdresse promo[TAILLE_MAX_PROMO];
    int nbEtudiants;
    string ville;

    // ... saisie de nbEtudiants et de promo (cf. exercice precedent) ...

    ville = get_string("Ville recherchee : ");

    printf("Etudiants de %s :\n", ville);
    (*\trou[{afficher\_ville(promo, nbEtudiants, ville);}]{13.5cm}*)

    printf("Moyenne a %s : %.2f\n", ville,
           (*\trou[{moyenne\_ville(promo, nbEtudiants, ville)}]{11.5cm}*));
    return 0;
}
\end{lstlisting}
\end{UPSTIexercice}
